\section{Introduction}\label{sec:introduction}

Let $\Lambda$ be a finite nonempty alphabet, and let
\[
 \Phi=(\varphi_i)_{i\in\Lambda},\qquad
 \varphi_i(x)=r_i x+t_i,\qquad t_i\in\R,\quad 0<|r_i|<1.
\]
The family is indexed: different symbols may specify the same map.
Given a probability vector $p=(p_i)_{i\in\Lambda}$ with $p_i>0$, its
self-similar measure is the unique Borel probability measure satisfying
\[
 \mu=\sum_{i\in\Lambda}p_i(\varphi_i)_*\mu.
\]
Equivalently, $\mu$ is the law of
$\lim_{n\to\infty}\varphi_{I_1}\circ\cdots\circ\varphi_{I_n}(0)$,
where the symbols $I_j$ are independent with law $p$.
The limit exists uniformly in the address: if
$r_{\max}=\max_i|r_i|$ and $T_0=\max_i|t_i|$, all coding limits
have absolute value at most $T_0/(1-r_{\max})$.

For a word $w=i_1\cdots i_n$, write
$\varphi_w=\varphi_{i_1}\circ\cdots\circ\varphi_{i_n}$ and
$p_w=p_{i_1}\cdots p_{i_n}$.  An \emph{exact overlap} is an equality
$\varphi_u=\varphi_v$ for distinct words of the same positive length.
Equality here means equality of the complete affine maps, including
both the signed slope and the translation.  We allow all such overlaps.

The relevant entropy counts maps rather than addresses.  Set
\[
 G_n=\varphi_{I_1}\circ\cdots\circ\varphi_{I_n},\qquad
 \P(G_n=g)=\sum_{w\in\Lambda^n:\,\varphi_w=g}p_w.
\]
All logarithms below have base two.  For a finite-valued random variable
$Z$, its Shannon entropy is $H(Z)=-\sum_z\P(Z=z)\log\P(Z=z)$,
with $0\log0=0$.  Define the random-walk entropy rate and the Lyapunov
exponent by
\begin{equation}\label{eq:rate-and-lyapunov}
 h=h_{\mathrm{RW}}(\Phi,p)
   :=\lim_{n\to\infty}\frac{H(G_n)}n
    =\inf_{n\ge1}\frac{H(G_n)}n,
 \qquad
 \chi=-\sum_i p_i\log|r_i|>0.
\end{equation}
Indeed, a length-$(n+m)$ map is a function of two independent maps
with laws $G_n,G_m$, so $H(G_{n+m})\le H(G_n)+H(G_m)$; subadditivity
gives the displayed limit and infimum.  Also
$0\le h\le H(p):=-\sum_i p_i\log p_i$.  The ratio $h/\chi$ is
independent of the common logarithm base.

We use the lower Hausdorff dimension of a measure:
\[
 \dim_H\nu=\inf\{\dim_H E:E\subset\R\text{ Borel},\ \nu(E)>0\}.
\]
Feng--Hu's exact-dimensionality theorem
\cite[Theorem~2.8]{FengHu2009} applies to finite self-similar systems
without separation.  Indeed, the maps preserve a sufficiently large
compact interval, extend to contracting smooth diffeomorphisms, and have
derivative norm and least singular value both equal to $|r_i|$.
The system is therefore conformal, and its Bernoulli coding law is
ergodic.  These observations also cover negative ratios and repeated
indexed maps.  Thus, for the measure
$\mu$ above, there is a constant $d$ such that
\begin{equation}\label{eq:exact-dimensionality}
 \lim_{s\downarrow0}\frac{\log\mu(B(x,s))}{\log s}=d
 \quad\text{for $\mu$-almost every $x$}.
\end{equation}
For an exact-dimensional measure this constant equals both $\dim_H\mu$
and $\inf\{\dim_H E:\mu(E)=1\}$.
To see the equivalence, restrict to countably many sets on which the
local bounds $s^{d+\varepsilon}\le\mu(B(x,s))\le s^{d-\varepsilon}$
hold uniformly for all sufficiently small $s$.  The lower mass bound
gives a full-measure union of sets of dimension at most $d+\varepsilon$
by a covering argument; the upper mass bound gives dimension at least
$d-\varepsilon$ for every positive-mass set.  Let
$\varepsilon\downarrow0$.  Consequently either Hausdorff
measure-dimension convention gives the same statement below.

\begin{theorem}\label{thm:main}
For every finite nonempty indexed family
$\varphi_i(x)=r_i x+t_i$ on $\R$ with $0<|r_i|<1$, and every
strictly positive probability vector $p$, its self-similar measure satisfies
\[
 \dim_H\mu_{\Phi,p}
 =\min\left\{1,\frac{h_{\mathrm{RW}}(\Phi,p)}{\chi(\Phi,p)}\right\}.
\]
No separation assumption is required; exact overlaps and repeated
generators are allowed.
\end{theorem}

This proves the entropy-rate dimension conjecture, formulated as
Conjecture~3 in Varj\'u's survey \cite{VarjuSurvey2026}.
When there are no exact overlaps, $H(G_n)=nH(p)$, and the theorem
gives the usual entropy-to-Lyapunov formula.  Section~\ref{sec:corollaries}
records this case, the attractor formula, and the homogeneous all-weight
consequence, including the three-map family
$\{\lambda x,\lambda x+1,\lambda x+t\}$.
The theorem includes critical and supercritical entropy
rates.  It asserts dimension, without an assertion of absolute continuity.

\subsection{Background and proof strategy}

The dimension of a self-similar set is classical when its pieces are
sufficiently separated.  Moran's construction \cite{Moran1946} and
Hutchinson's iterated-function-system framework \cite{Hutchinson1981}
give the similarity-dimension formula under the open set condition.
Hutchinson also established the invariant-measure construction used
above \cite[Theorem~4.4(1)]{Hutchinson1981}.
Removing separation leads to the exact-overlap problem: can dimension
fall below its natural upper bound when distinct words never define
the same map?  The set version is associated with Simon
\cite{Simon1996}; see also the formulation and historical discussion
in \cite[Section~1]{VarjuSurvey2026}.
The entropy-rate formulation addresses the additional loss of
information when exact overlaps do occur.

Bernoulli convolutions provide an early and influential model for this
question.  These are the laws of $\sum_{j\ge0}\varepsilon_j\lambda^j$,
where $0<\lambda<1$ and the $\varepsilon_j$ are independent
equiprobable signs.
Erd\H{o}s proved singularity for reciprocal Pisot parameters
in $(1/2,1)$ \cite{Erdos1939}.  A Pisot number is a real algebraic
integer greater than one whose other conjugates have modulus less
than one.  Solomyak proved absolute continuity with an
$L^2$ density for almost every parameter in $(1/2,1)$
\cite{Solomyak1995}; Peres--Solomyak later gave a simpler proof
\cite{PeresSolomyak1996}.
Garsia introduced the discrete entropy rate in this setting and
related it to singularity \cite[Theorem~1.2]{Garsia1963}.
These results illustrate the sensitivity to overlaps and arithmetic;
absolute continuity and Hausdorff dimension are distinct questions.

Hochman showed that $\dim_H\mu<\min\{1,H(p)/\chi\}$
forces superexponential concentration of cylinder maps
\cite[Theorem~1.1]{Hochman2014}, using inverse theorems for entropy
growth under convolution.  For unequal ratios his entropy formulation
already keeps track of the full affine map, including its slope
\cite[Theorem~1.4]{Hochman2014}.
The entropy-rate formula under weak exponential separation follows
from his work; the version allowing exact collisions is proved in
B\'ar\'any--Verma \cite[Theorem~3.5]{BaranyVerma2026}.
For Bernoulli convolutions, Breuillard--Varj\'u developed approximation
by algebraic parameters with controlled entropy
\cite{BreuillardVarju2019}.  Building on these and earlier entropy
methods, Varj\'u proved full dimension for every transcendental parameter
in $(1/2,1)$ \cite[Theorem~3]{VarjuTranscendental2019}.
Rapaport proved the no-exact-overlap formula for algebraic
contraction ratios and arbitrary real translations, allowing signed
unequal ratios \cite[Theorem~2]{Rapaport2022}.

Rapaport--Varj\'u extended entropy and approximation methods to
homogeneous three-map systems \cite{RapaportVarju2024}.
For equal weights, their results include an exceptional parameter set
of Hausdorff dimension zero and the no-overlap formula when $\lambda>2^{-2/3}$
\cite[Corollary~1.4 and Theorem~1.8]{RapaportVarju2024}.
Rapaport--Varj\'u proved the entropy-rate formula, including exact overlaps,
for rational translations and a positive common ratio
\cite[Theorem~A.1]{RapaportVarju2024}.  Feng--Feng proved the
no-exact-overlap formula for algebraic translations and a signed common ratio
\cite[Theorem~1.2]{FengFeng2025}.
The distinction in Theorem~\ref{thm:main} is that the entropy rate of
the full affine-map walk is sharp even when exact collisions occur,
with no arithmetic or separation assumption.

Absence of exact overlaps does not itself supply a quantitative
separation estimate.  Baker \cite[Theorem~1.3]{Baker2021} and,
independently, B\'ar\'any--K\"aenm\"aki
\cite[Theorem~2.1]{BaranyKaenmaki2021} constructed systems without
exact overlaps whose distinct cylinders approach one another at
arbitrarily prescribed superexponential rates.
Thus the entropy carried by extremely close maps must be handled
without bounding the smallest positive separation.

Our proof combines the entropy retained by this map walk with a
finite-law estimate that detects information across arbitrarily fine scales.
It uses translation averaging, developed by Wang
\cite[Section~4.1]{Wang2011}, and the entropy calculus and fair-pair
arguments of Varj\'u
\cite[Section~2.2, Propositions~20--21, and Lemma~22]{VarjuAbsolute2019}.
Uniform nonsaturation has an antecedent for Bernoulli convolutions in
Breuillard--Varj\'u \cite[Lemma~13]{BreuillardVarju2019}.
We prove the required entropy estimates here; the only external
dimension theorem used in the proof is exact dimensionality.

The argument has three steps.
First, for a finite real law $\nu$, compare its exact entropy with
the entropy of its grid cell at a fixed scale.  The difference is
information hidden inside those cells.  A bound for the sum of several
independent copies forces this information to be witnessed by fair
(equal-weight) two-point laws at smaller scales.  The estimate is independent of the smallest
distance between support points.

Second, group the coding symbols into blocks and condition on their
symbol counts.  This is the block-type disintegration of
Galicer--Saglietti--Shmerkin--Yavicoli
\cite[Section~6.4 and Lemma~6.6]{GalicerSagliettiShmerkinYavicoli2016},
also developed by Saglietti--Shmerkin--Solomyak
\cite[Lemma~6.2]{SagliettiShmerkinSolomyak2018};
K\"aenm\"aki--Orponen \cite[Section~2.2]{KaenmakiOrponen2023}
explicitly allow repeated indexed maps in this construction.
Here the counts fix each block's signed contraction.
For a sequence of such blocks, its translation then identifies its
complete composed map.  Conditioning costs only the entropy of the
counts, so the map-entropy rate gives a lower bound for the remaining
translation entropy even when addresses collide.  The same fixed
contractions give a small sumset bound.  A hypothetical dimension
deficit therefore yields a linear amount of fair-pair mass across a
finite band of very fine scales.

Third, each pair produces a positive entropy gain against the
unobserved tail.  The bands may end at arbitrarily large depths.
We choose successive block lengths only after the preceding endpoints
are known.  At each observation scale their retained blocks are
disjoint, so their gains telescope to a bounded total.  Averaging over
observation scales gives a fixed positive contribution from every
band, a contradiction as the number of bands grows.

Section~\ref{sec:entropy} proves the entropy calculus, nonsaturation,
and two-point gain.  Section~\ref{sec:pairs} proves the finite-law
estimate.  Section~\ref{sec:types} applies it to conditional map laws,
and Section~\ref{sec:windows} arranges the disjoint gains and completes
the proof.  Section~\ref{sec:corollaries} derives the no-overlap measure
and attractor formulas, followed by the homogeneous and three-map
specializations.  No quantitative bound on the fine-scale endpoints is needed.

\subsection{The entropy-rate upper bound}\label{sec:upper}

We prove the upper bound before turning to the contradiction argument.
Fix a positive integer $n$.  Use the finite set of distinct maps in
the support of $G_n$ as a new alphabet, giving each map $g$ its exact
probability $q_g=\P(G_n=g)>0$.  Grouping independent coding symbols
into $n$-blocks shows that this new system has the same measure $\mu$.
Its one-symbol entropy is $H(G_n)$ and its mean contraction depth is
$n\chi$, since coincident maps have the same absolute slope.

For $0<\varepsilon<n\chi$, the strong law shows that almost every
address in this block system eventually has, at every length $j$,
a prefix of probability at least $2^{-j(H(G_n)+\varepsilon)}$ and
absolute contraction at most $2^{-j(n\chi-\varepsilon)}$.
There are at most $2^{j(H(G_n)+\varepsilon)}$ such prefixes.
For each fixed eventual starting length $J$, let $E_J$ be the closed
set of addresses satisfying both conditions for every $j\ge J$.
Its compact coding image is covered by intervals of diameter at most
$C2^{-j(n\chi-\varepsilon)}$, where $C$ bounds the attractor diameter.
The total $a$th power of these diameters tends to zero when
$a>(H(G_n)+\varepsilon)/(n\chi-\varepsilon)$.
The countable union over eventual starting lengths is a full-measure
Borel set with this dimension bound.  The case of zero attractor
diameter is immediate.  Letting $\varepsilon\downarrow0$, then taking
the infimum over $n$, gives
\begin{equation}\label{eq:upper-bound}
 d\le\min\{1,h/\chi\}.
\end{equation}
In particular, the conclusion of Theorem~\ref{thm:main} already holds
when $h=0$, including a single-map system.  To prove equality in all
other cases, it suffices to rule out
\begin{equation}\label{eq:dimension-loss}
 d<1,\qquad d\chi<h.
\end{equation}
These inequalities are the contradiction hypothesis in
Sections~\ref{sec:types} and~\ref{sec:windows}.
